% =====================================================================
% Lean–Julia Integration Notes v2.0
% Constellation-Manifold-Toolkit — crossanalysis/devnotes
% =====================================================================

\section{Lean–Julia Integration Notes v2.0}
\label{sec:LeanJuliaIntegrationV2}

This document provides updated integration guidance for Lean formal modules
and Julia runtime modules in light of the Enhanced Conceptual Collision
Tracking framework, including Φ7 curvature envelopes, Φ8 drift semigroup
evolution, Φ9 invariant extraction, collapse-surface diagnostics, and
hallucination-guard constraints.

% ---------------------------------------------------------------------
\subsection{Lean Module Population Strategy (v2.0)}

Lean modules must be populated before Julia runtime modules. Lean provides:

\begin{itemize}
  \item formal invariants for conceptual manifolds,
  \item curvature-envelope safety proofs,
  \item drift-semigroup semigroup properties,
  \item invariant-extraction correctness,
  \item deformation-tensor safety,
  \item collision-avoidance theorems.
\end{itemize}

Lean-Hallucination-Guard v2.0 must be applied to every definition:

\[
\#guard\; ConceptualManifold,\quad
\#guard\; CurvatureEnvelope,\quad
\#guard\; DriftSemigroup,\quad
\#guard\; InvariantExtraction,\quad
\#guard\; DeformationTensor,\quad
\#guard\; CollisionInvariant.
\]

This prevents implicit type inference, operator drift, and altitude contamination.

% ---------------------------------------------------------------------
\subsection{Lean Module Structure (Updated)}

Lean modules required for v2.0:

\begin{verbatim}
crossanalysis/lean/
  ConceptualManifold.lean
  CurvatureEnvelope.lean
  CollapseSurface.lean
  DriftSemigroup.lean
  InvariantExtraction.lean
  SoftDeformation.lean
  PhaseShiftOperator.lean
  EnvelopeSafety.lean
  CollisionInvariant.lean
\end{verbatim}

New modules in v2.0:

\begin{itemize}
  \item CollapseSurface.lean
  \item DriftSemigroup.lean (Φ8)
  \item InvariantExtraction.lean (Φ9)
\end{itemize}

These correspond to the enhanced conceptual collision tracking spine.

% ---------------------------------------------------------------------
\subsection{Julia Runtime Integration (v2.0)}

Julia modules must mirror Lean modules:

\begin{verbatim}
crossanalysis/julia/
  ConceptualManifold.jl
  CurvatureEnvelope.jl
  CollapseSurface.jl
  DriftSemigroup.jl
  InvariantExtraction.jl
  SoftDeformation.jl
  PhaseShiftOperator.jl
  EnvelopeSafety.jl
  CollisionRuntime.jl
\end{verbatim}

Julia must enforce Lean-verified invariants:

\[
\text{assert}(R'^{(c)} < R_{\text{crit}}),\quad
\text{assert}(D_c(t) < D_{\text{crit}}),\quad
\text{assert}(\det(g'^{(c)}) \neq 0).
\]

Julia must not infer operator signatures or tensor ranks.

Julia must not introduce new invariants.

Julia must not mix altitudes.

% ---------------------------------------------------------------------
\subsection{Hallucination Guard v2.0 Requirements}

Lean-Hallucination-Guard v2.0 introduces:

\begin{itemize}
  \item operator-rank verification,
  \item tensor-shape verification,
  \item curvature-envelope domain checks,
  \item drift-semigroup closure checks,
  \item invariant-extraction stability checks,
  \item deformation-tensor smoothness checks.
\end{itemize}

Every Lean definition must be guarded.

Every Julia runtime call must reference a Lean-verified invariant.

% ---------------------------------------------------------------------
\subsection{Enhanced Conceptual Collision Tracking Integration}

Enhanced collision tracking requires:

\[
\Phi_7^{(c)}\ \text{(curvature envelope)},\quad
\Phi_{8,c}^t\ \text{(drift semigroup)},\quad
\Phi_{9,c}\ \text{(invariant extraction)}.
\]

Lean must prove:

\[
\lambda_{\min}(H_c(x)) < 0 \Rightarrow x \in \Sigma_c,
\]
\[
D_c(t) > D_{\text{crit}} \Rightarrow \text{collision precursor},
\]
\[
I_k(\rho_1) \approx I_k(\rho_2) \Rightarrow \text{invariant convergence}.
\]

Julia must implement:

\begin{itemize}
  \item curvature-envelope evaluation,
  \item drift-semigroup evolution,
  \item invariant-extraction runtime diagnostics,
  \item collapse-surface scanning,
  \item phase-shift steering,
  \item deformation-tensor modulation.
\end{itemize}

% ---------------------------------------------------------------------
\subsection{Altitude Safety}

Lean modules correspond to A0–A2 (reversible geometry).

Julia modules correspond to A3–A4 (runtime physics and diagnostics).

No module may cross altitudes.

% =====================================================================
% End Lean–Julia Integration Notes v2.0
% =====================================================================

---
Provenance:
Developed as part of the Constellation-Manifold-Toolkit cross-analysis expansion.
Provides updated Lean–Julia integration guidance for Enhanced Conceptual Collision
Tracking, incorporating Φ7 curvature envelopes, Φ8 drift semigroup evolution,
Φ9 invariant extraction, collapse-surface diagnostics, and Hallucination Guard v2.0.
Altitude-safe and operator-bounded. Not part of TISD governance, Anima-Vectorium
runtime, or Vectorium A7–A9. Serves exclusively as technical developer documentation.

Author:
Borealis S. Hedling — Dublin, Ireland.

Version:
CMT-LJI-v2.0 — 27 October 2026, 21:14 IST
---
